\section{变分自编码器（VAE）回顾}

\subsection{VAE 的核心思想}

VAE 的目标是学习解码器（decoder）$p_\theta(x|z)$ 以最大化数据的边际似然 $p_\theta(x)$。然而，直接计算边际似然需要对所有可能的 $z$ 做积分，这在计算上是不可行的：

$$p_\theta(x) = \int p_\theta(x|z) p(z) \, dz$$

虽然我们可以通过蒙特卡洛估计来近似该积分，但这需要采样大量的 $z$，计算代价极高。

\begin{warningbox}{边际似然不可直接计算}
通过贝叶斯公式 $p_\theta(x) = \frac{p_\theta(x|z) p(z)}{p_\theta(z|x)}$，计算边际似然还需要知道后验分布 $p_\theta(z|x)$，但后验本身也是未知的。这构成了一个``鸡生蛋、蛋生鸡''的困境，正是 VAE 引入变分后验 $q_\phi(z|x)$ 的根本原因。
\end{warningbox}

\subsection{ELBO：证据下界}

由于无法直接最大化 $\log p_\theta(x)$，VAE 转而最大化其下界——\textbf{证据下界}（Evidence Lower Bound, ELBO）。

ELBO 的推导基于 Jensen 不等式。对于任意分布 $q_\phi(z|x)$：

$$\log p_\theta(x) = \log \int p_\theta(x|z) p(z) \, dz$$

在积分中引入 $q_\phi(z|x)$：

$$= \log \int \frac{p_\theta(x|z) p(z)}{q_\phi(z|x)} q_\phi(z|x) \, dz = \log \mathbb{E}_{q_\phi(z|x)} \left[ \frac{p_\theta(x|z) p(z)}{q_\phi(z|x)} \right]$$

由 Jensen 不等式（$\log$ 是凹函数）：

$$\geq \mathbb{E}_{q_\phi(z|x)} \left[ \log \frac{p_\theta(x|z) p(z)}{q_\phi(z|x)} \right] = \text{ELBO}$$

\begin{importantbox}{ELBO 的分解}
ELBO 可以分解为两项：
$$\text{ELBO} = \underbrace{\mathbb{E}_{q_\phi(z|x)}[\log p_\theta(x|z)]}_{\text{重构项（Reconstruction Term）}} - \underbrace{D_{\text{KL}}(q_\phi(z|x) \| p(z))}_{\text{KL 正则项}}$$

\begin{itemize}
    \item \textbf{重构项}：鼓励解码器从潜变量 $z$ 准确重构输入 $x$，本质上等价于自编码器的重构损失。
    \item \textbf{KL 正则项}：鼓励变分后验 $q_\phi(z|x)$ 尽可能接近先验 $p(z) = \mathcal{N}(0, I)$。
\end{itemize}
\end{importantbox}

\subsection{VAE 的训练过程}

VAE 同时训练编码器和解码器：

\begin{itemize}
    \item \textbf{编码器}（Encoder）：参数化变分后验 $q_\phi(z|x) = \mathcal{N}(\mu_\phi(x), \sigma^2_\phi(x) I)$，将数据映射到潜变量空间的均值和方差。
    \item \textbf{解码器}（Decoder）：参数化似然 $p_\theta(x|z) = \mathcal{N}(\mu_\theta(z), \sigma^2 I)$，将潜变量映射回数据空间。
\end{itemize}

\begin{knowledgebox}{重参数化技巧（Reparameterization Trick）}
采样操作 $z \sim q_\phi(z|x)$ 本身不可微分，但通过重参数化可以绕过这个问题：

$$z = \mu_\phi(x) + \sigma_\phi(x) \odot \epsilon, \quad \epsilon \sim \mathcal{N}(0, I)$$

这样，$z$ 关于 $\mu_\phi$ 和 $\sigma_\phi$ 是可微的，梯度可以正常反向传播到编码器参数 $\phi$。
\end{knowledgebox}

在生成阶段，只需使用解码器：从 $p(z) = \mathcal{N}(0, I)$ 采样 $z$，然后通过解码器 $p_\theta(x|z)$ 生成 $x$。编码器仅在训练阶段用于学习更好的似然分布。

\subsection{VAE 的两个高斯 KL 散度}

在 VAE 的训练中，需要计算两个各向同性高斯分布之间的 KL 散度。对于 $p = \mathcal{N}(\mu, \sigma^2 I)$ 和 $q = \mathcal{N}(0, I)$：

$$D_{\text{KL}}(p \| q) = \frac{1}{2} \left( \|\mu\|^2 + \text{tr}(\sigma^2 I) - d - \log \det(\sigma^2 I) \right)$$

$$= \frac{1}{2} \sum_{i=1}^{d} \left( \mu_i^2 + \sigma_i^2 - 1 - \log \sigma_i^2 \right)$$

其中 $d$ 是潜变量的维度。

\subsection{ELBO 的紧致性与 VAE 的局限}

\begin{importantbox}{ELBO 紧致条件}
ELBO 等于真实的对数似然，当且仅当变分后验与真实后验完全一致：
$$\text{ELBO} = \log p_\theta(x) \iff q_\phi(z|x) = p_\theta(z|x)$$

两者之间的差距恰好是：$\log p_\theta(x) - \text{ELBO} = D_{\text{KL}}(q_\phi(z|x) \| p_\theta(z|x))$，这个差距被称为\textbf{松弛度}（tightness gap）。
\end{importantbox}

\begin{warningbox}{VAE 的根本局限}
VAE 将似然分布 $p_\theta(x|z)$ 和变分后验 $q_\phi(z|x)$ 都建模为高斯分布。然而，对于复杂的数据分布（如自然图像），\textbf{同时使似然和后验都为高斯分布在数学上通常是不可能的}。这种过度简化导致 VAE 生成的图像往往较为模糊。
\end{warningbox}

VAE 的这一局限直接激发了对更复杂生成模型的探索——特别是通过引入\textbf{层级化结构}来增加模型的表达能力。

\subsection{本章小结}

VAE 通过 ELBO 绕过了边际似然不可计算的问题，但其高斯假设过于简单，限制了生成质量。为了在保持计算可行性的同时提高表达能力，下一步是引入层级化的 VAE 结构。


